\scanpage{31}
\mspage{15}
\unsure{lifted|to lift} to derived hypotheses arithmetico-geometric \inserted{in nature}, one can drop the anabelian aspects, to play to the full on the abelian aspects -- but in an arithmetic context. Instead of having \inserted{and using} only, as in corollary 1 a), hypotheses that $H_1(\overline f)=H_1(\overline g):H_1(\overline X)\longrightarrow H_1(\overline Y)$, one has a finer hypothesis, on the equality of two homomorphisms \uline{of extensions} of $E_{\overline K/K}$ by $H_1(\overline X,\widehat{\mathbb Z}_\ell)$ respectively $H_1(\overline Y,\widehat{\mathbb Z}_\ell)$. $\forall\ell\in\mathbb P$.

Proof of corollary 3.
It suffices to see that $f(x)=g(x)$ for every $x\in X(K)$ (for, after possibly \deleted{extending} \inserted{making a finite extension} of the base field, this remains true \deleted{after any field extension} \inserted{for $x\in X(\overline K)$}). \deleted{\unsure{do}}

Now a point defines a \inserted{conjugacy} class of \deleted{on the other hand} splittings of $E_{X/K}$ over $E_{\overline K/K}$, and a fortiori of the associated extension with abelian kernel -- the hypothesis implies that $f(x)$ and $g(x)$ define the $\unsure{=|\overline m}$ class of
\margin{30}
